% Generated draft: requires mathematical and source review.
\documentclass[11pt]{article}
\usepackage[margin=25mm]{geometry}
\usepackage{fontspec}
\setmainfont{Latin Modern Roman}
\usepackage{amsmath,amssymb,mathtools,unicode-math}
\setmathfont{Latin Modern Math}
\usepackage{enumitem,needspace,xurl,hyperref}
\setlength{\parindent}{0pt}
\setlength{\parskip}{5pt}
\setlength{\emergencystretch}{3em}
\newcommand{\N}{\mathbb{N}}
\newcommand{\Z}{\mathbb{Z}}
\newcommand{\Q}{\mathbb{Q}}
\newcommand{\R}{\mathbb{R}}
\newcommand{\C}{\mathbb{C}}
\newcommand{\F}{\mathbb{F}}
\newcommand{\A}{\mathbb{A}}
\newcommand{\PP}{\mathbb{P}}
\newcommand{\E}{\mathbb{E}}
\newcommand{\eps}{\varepsilon}
\newcommand{\dd}{\,\mathrm{d}}
\begin{document}
\section*{UnsolvedMath ID: 30000846}
\textbf{Source problem code:} catalogue entry
\subsection{Definitions and mathematical statement}
Let \(G\) be a finite group, let \(p\) be a prime, and let \(k\) be a field of characteristic \(p\). The \(p\)-orbit category \(O_p(G)\) has objects \(G/H\), where \(H\leq G\) is a \(p\)-subgroup, and morphisms the \(G\)-equivariant maps. Its \(k\)-linearization allows the study of contravariant orbit-category modules, equivalently functors from \(O_p(G)^{\mathrm{op}}\) to \(k\)-vector spaces.

The source asks whether the following equivalence holds for every oriented dimension function \(f\) in the sense of the report:
\[
\exists\text{ a }kO_p(G)^{\mathrm{op}}\text{-sphere }C
\text{ realizing }f
\quad\Longleftrightarrow\quad
f\text{ satisfies the Borel--Smith conditions}.
\]
A dimension function records the sphere dimension at each subgroup object and is invariant under conjugation. In the usual homological interpretation, evaluation at \(G/H\) has reduced homology concentrated in the degree prescribed by \(f(H)\). An orientation requirement typically controls the action of the Weyl group \(N_G(H)/H\) on the top homology, but the exact definition of ``oriented dimension function'' in this problem is not supplied and is not replaced here by a guessed parity or action condition.

The extraction also omits the precise meaning of an orbit-category sphere: whether it is a finite free or projective chain complex, what additional realization axioms it satisfies, the permitted dimension range and degree shift, and the full list of Borel--Smith relations. These definitions are necessary to turn the displayed equivalence into a self-contained theorem statement. In particular, realization by an algebraic orbit-category sphere must not be silently replaced by realization by a genuine sphere with a \(G\)-action, a representation sphere, or a spherical fibration. The source asks an exact realization equivalence; it does not state merely realization after taking joins or multiplying the dimension function.
\subsection{Short English statement}
Are the Borel--Smith conditions exactly the criterion for realizing every oriented dimension function by a sphere over the \(p\)-orbit category? The precise algebraic sphere and orientation definitions must be recovered from the report.
\subsection{Sources}
\begin{itemize}
\item Local UnsolvedMath registry, permanent ID \(30000846\), source code catalogue entry, ``Realization of Oriented Dimension Functions by Equivariant Spheres.'' No individual proposer is supplied.
\item \emph{Homotopy Theory} (2008), Oberwolfach Reports, DOI 10.4171/owr/2007/45, pp. 2671--2726: \url{https://doi.org/10.4171/owr/2007/45}.
\item Supplied reference, ``Okay--Yalcin, Dimension functions for spherical fibrations'': \url{https://arxiv.org/abs/1710.07315}.
\item Supplied literature triage, checked 2026-08-21: status ``partially solved.'' It reports necessity of Borel--Smith conditions, sufficiency for classes including many small groups and groups with normal Sylow subgroups, and further necessary conditions after joins in later work. The exact special-case hypotheses are absent. This dated assessment derives from supplied OWR curation; time-sensitive status and cited claims require independent verification.
\end{itemize}
\end{document}
